\section{Reconstructed small normal budgets for R1 unbounded-$q$}
\label{sec:reconstructed-r1-qo-small-normal-budgets}

\paragraph{Scope and provenance.}
This is a fresh proof from the actual compact source RSR1--RSR9 and
the exact circuit identity RSR23. It replaces the needed mathematics
of the lost NS helper; it is not an assertion of original file bytes.
Use all notation of \texttt{r1\_qo\_actual\_source\_and\_scalar.tex}.
The testing variables $u,v>0$ are arbitrary independent infinitesimals,
and $z>0$ is arbitrary. Put
\[
 b=v-uz/j=\sqrt y\,B_*,\qquad
 h_0=\mu+\tfrac12,\qquad
 M_e=\frac{nh_0j^2}{n+h_0},\qquad J_c=M_ey.
 \tag{NSR1}
\]
The two rows
\[
       p v^2b,\qquad \sqrt T\,\mu u v b
 \tag{NSR2}
\]
can be given any fixed real positive square coefficient at
infinitesimal cost from $f_c\Psi+J_cB_*^2$. The same holds after
multiplying either row by any bounded separated coefficient.
No coefficient in the budget depends on $z$.

\paragraph{The infinitesimal-$\eta_*$ source case.}
By RSR9, $n>0$ and
$n\sim p/[2(j-p)(j+1)]$. The exact identity
\[
 \frac{p^2}{M_e}
 =\frac{p^2}{h_0j^2}+\frac{p^2}{nj^2}
 \tag{NSR3}
\]
is bounded: the first term is bounded since $p/j<1$; the second
is comparable to $p(j-p)(j+1)/j^2$. If $j\leq1$ this last expression
is at most $2p/j$, and if $j\geq1$ it is at most $2p$.
Thus the direct costs of the squared rows in NSR2, relative to
$J_cB_*^2=M_eb^2$, are
\[
 \frac{p^2}{M_e}v^4\in\mathrm I,\qquad
 \frac{T\mu^2}{M_e}u^2v^2\in\mathrm I.
 \tag{NSR4}
\]
Here $\mu/p$ is bounded by RSR6--RSR7. No use of the circuit is
needed in this case, and no positive lower bound on $j$ is imposed.

\paragraph{The positive-unit-$\eta_*$ source case.}
The remaining possibility is that $\eta_*$ has a positive real lower
bound, since RSR9 already proves it bounded. The exact relation
\[
 \frac{p^2}{j^2B}=\frac{p+p/j}{\eta_*}
\]
and RSR16--RSR17 imply $p^2y/f_c$ bounded. Apply
\[
 (1+V^2)\Psi-V^2B_*^2
       =Q_*^2+R_*^2+V^2(1-UVZ_*)^2.
 \tag{NSR5}
\]
The first row in NSR2 costs exactly
\[
 \frac{p^2y}{f_c}v^2(y+v^2)\in\mathrm I
\]
on $f_c\Psi$. Indeed the difference between that multiple of the
circuit packet and $(pv^2b)^2$ is $p^2y^3V^2$ times the SOS on
the right of NSR5. The second row costs exactly
\[
 \frac{T\mu^2y}{f_c}u^2(y+v^2)\in\mathrm I;
\]
its difference is $T\mu^2xy^2U^2$ times the same SOS. Every
multiplier is a product of nonnegative pure-source and pure-testing
factors. These formulas include literal $n=0$, since they never
divide by the current mass. The two source cases exhaust the source.

\paragraph{Use in the full normal conductor.}
The first row pays the $(p/2)v^2b$ term in the complete half-$N$
derivative. The second pays the $\mu uvb$ term in the half-$M$
derivative after its source normalization by
$\sqrt{TH(1-H)}$: the omitted factor $\sqrt{H(1-H)}$ is bounded.
This is only a square-budget lemma. The exact full-normal identities
and the retained hard block are required separately to infer a
full-parent local certificate.
